% ====================================================================
% ФАЙЛ: tasks/01_mechanics/newton_second_law_bank.tex
% ТЕМА: ВТОРОЙ ЗАКОН НЬЮТОНА (ДИНАМИКА И КРИВОЛИНЕЙНОЕ ДВИЖЕНИЕ)
% ПАЧКА 1: ВАРИАНТЫ 1–10
% ====================================================================

% === ЗАДАЧА 1 ===
% Тег: #МЕХАНИКА #КИМ_02 #ВАРИАНТ_01
\begin{taskbasic}[code={\label{task:newton2_v1_n2}}]
\textbf{Задача 1.} В инерциальной системе отсчёта сила $F$ сообщает телу массой $m$ ускорение $a$. Во сколько раз нужно увеличить массу тела, чтобы вдвое большая сила сообщала ему в той же системе отсчёта в $4$ раза меньшее ускорение?

\kim{2}
\end{taskbasic}
\answer{8}

% === ЗАДАЧА 2 ===
% Тег: #МЕХАНИКА #КИМ_05 #ВАРИАНТ_01
\begin{taskbasic}[code={\label{task:newton2_v1_n5}}]
\textbf{Задача 2.} Грузовик массой $10$~т проезжает верхнюю точку выпуклого моста, радиус кривизны которого равен $80$~м, двигаясь равномерно со скоростью $72$~км/ч. Из приведённого ниже списка выберите все верные утверждения, характеризующие движение грузовика.

1) Сила, с которой мост действует на грузовик, больше $40$~кН и направлена вертикально вверх.
2) Сумма сил, действующих на грузовик, направлена вертикально вверх и перпендикулярна его скорости.
3) Сила, с которой грузовик действует на мост, направлена вертикально вниз и равна $100$~кН.
4) Сила тяжести, действующая на грузовик, равна $100$~кН.
5) Центростремительное ускорение грузовика равен $5$~$м/с^2$.

\kim{5}
\end{taskbasic}
\answer{145}

% === ЗАДАЧА 3 ===
% Тег: #МЕХАНИКА #КИМ_02 #ВАРИАНТ_02
\begin{taskbasic}[code={\label{task:newton2_v2_n2}}]
\textbf{Задача 3.} В инерциальной системе отсчёта сила $F$ сообщает телу массой $m$ ускорение $a$. Во сколько раз нужно уменьшить массу тела, чтобы втрое меньшая сила сообщала ему в той же системе отсчёта в $2$ раза большее ускорение?

\kim{2}
\end{taskbasic}
\answer{6}

% === ЗАДАЧА 4 ===
% Тег: #МЕХАНИКА #КИМ_05 #ВАРИАНТ_03
\begin{taskbasic}[code={\label{task:newton2_v3_n5}}]
\textbf{Задача 4.} В лабораторной работе изучали движение небольшого бруска массой $200$~г по горизонтальной шероховатой поверхности под действием горизонтальной постоянной силы, равной по модулю $800$~мН. Зависимость модуля скорости бруска от времени приведена в таблице.
\begin{center}
\begin{tabular}{|c|c|c|c|c|c|c|c|}
\hline
Время $t$, с & 0 & 0,1 & 0,2 & 0,3 & 0,4 & 0,5 & 0,6 \\ \hline
Скорость $v$, м/с & 0 & 0,2 & 0,4 & 0,6 & 0,8 & 1,0 & 1,2 \\ \hline
\end{tabular}
\end{center}
Выберите все верные утверждения на основании анализа представленной таблицы.

1) Модуль ускорения бруска равен $2$~$м/с^2$.
2) Коэффициент трения бруска о поверхность $\mu = 0{,}3$.
3) Брусок движется равноускоренно.
4) В момент времени $0{,}3$~с кинетическая энергия бруска равна $36$~мДж.
5) Модуль равнодействующей сил, действующих на брусок, равен $0{,}8$~Н.

\kim{5}
\end{taskbasic}
\answer{134}

% === ЗАДАЧА 5 ===
% Тег: #МЕХАНИКА #КИМ_05 #ВАРИАНТ_04
\begin{taskbasic}[code={\label{task:newton2_v4_n5}}]
\textbf{Задача 5.} В лабораторной работе изучали движение небольшого бруска массой $200$~г по горизонтальной шероховатой поверхности под действием горизонтальной постоянной силы, равной по модулю $800$~мН. Зависимость модуля скорости бруска от времени приведена в таблице.
\begin{center}
\begin{tabular}{|c|c|c|c|c|c|c|c|}
\hline
Время $t$, с & 0 & 0,1 & 0,2 & 0,3 & 0,4 & 0,5 & 0,6 \\ \hline
Скорость $v$, м/с & 0 & 0,2 & 0,4 & 0,6 & 0,8 & 1,0 & 1,2 \\ \hline
\end{tabular}
\end{center}
Выберите все верные утверждения на основании анализа представленной таблицы.

1) Модуль ускорения бруска равен $0{,}4$~$м/с^2$.
2) Коэффициент трения бруска о поверхность $\mu = 0{,}2$.
3) Брусок движется равномерно.
4) В момент времени $0{,}3$~с кинетическая энергия бруска равна $18$~мДж.
5) Модуль равнодействующей сил, действующих на брусок, равен $0{,}4$~Н.

\kim{5}
\end{taskbasic}
\answer{25}

% === ЗАДАЧА 6 ===
% Тег: #МЕХАНИКА #КИМ_05 #ВАРИАНТ_05
\begin{taskbasic}[code={\label{task:newton2_v5_n5}}]
\textbf{Задача 6.} Два тела, каждое массой $2$~кг, движутся по одной прямой, вдоль которой направлена ось $Ox$. На рисунке приведены графики зависимости проекций $v_x$ их скоростей от времени $t$.
\begin{center}
\begin{tikzpicture}[x=0.4cm, y=0.35cm, every node/.style={font=\footnotesize}]
    \draw[xstep=2, ystep=1, gray!30, thin] (0,0) grid (18,7);
    \draw[-{Stealth[scale=0.8]}, black, thick] (0,0) -- (19.5,0) node[right] {$t, \text{ с}$};
    \draw[-{Stealth[scale=0.8]}, black, thick] (0,0) -- (0,7.5) node[above] {$v_x, \text{ м/с}$};
    \foreach \y in {2,4,6} \node[left] at (0,\y) {\y};
    \foreach \x in {4,8,12,16} \node[below] at (\x,0) {\x};
    \node[below left] at (0,0) {0};
    \draw[ultra thick, brandPrimary] (0,7) -- (18,2.8) node[right, black] {1};
    \draw[ultra thick, brandAccent] (0,0.5) -- (18,5.7) node[right, black] {2};
\end{tikzpicture}
\end{center}
Из приведённого ниже списка выберите все верные утверждения о движении тел.

1) Равнодействующая сил, действующих на тело 1, монотонно убывает в течение всего времени наблюдения.
2) В промежутке времени от $0$ до $14$~с тела двигались навстречу друг другу.
3) Кинетическая энергия тела 2 за промежуток времени от $6$ до $14$~с увеличилась в $4$ раза.
4) Путь, пройденный телом 1 за промежуток времени от $0$ до $18$~с, больше пути, пройденного телом 2 за тот же промежуток времени.
5) Модуль импульса тела 2 в момент времени $t=10$~с равен $10$~\text{кг}$\cdot$\text{м/с}.

\kim{5}
\end{taskbasic}
\answer{34}

% === ЗАДАЧА 7 ===
% Тег: #МЕХАНИКА #КИМ_05 #ВАРИАНТ_06
\begin{taskbasic}[code={\label{task:newton2_v6_n5}}]
\textbf{Задача 7.} Два тела, каждое массой $2$~кг, движутся по одной прямой, вдоль которой направлена ось $Ox$. На рисунке приведены графики зависимости проекций $v_x$ их скоростей от времени $t$.
\begin{center}
\begin{tikzpicture}[x=0.4cm, y=0.35cm, every node/.style={font=\footnotesize}]
    \draw[xstep=2, ystep=1, gray!30, thin] (0,0) grid (18,7);
    \draw[-{Stealth[scale=0.8]}, black, thick] (0,0) -- (19.5,0) node[right] {$t, \text{ с}$};
    \draw[-{Stealth[scale=0.8]}, black, thick] (0,0) -- (0,7.5) node[above] {$v_x, \text{ м/с}$};
    \foreach \y in {2,4,6} \node[left] at (0,\y) {\y};
    \foreach \x in {4,8,12,16} \node[below] at (\x,0) {\x};
    \node[below left] at (0,0) {0};
    \draw[ultra thick, brandPrimary] (0,7) -- (18,2.8) node[right, black] {1};
    \draw[ultra thick, brandAccent] (0,0.5) -- (18,5.7) node[right, black] {2};
\end{tikzpicture}
\end{center}
Из приведённого ниже списка выберите все верные утверждения о движении тел.

1) Равнодействующая сил, действующих на тело 1, остаётся в течение всего времени наблюдения постоянной.
2) В промежутке времени от $0$ до $14$~с тела двигались в одном направлении.
3) Кинетическая энергия тела 2 за промежуток времени от $6$ до $14$~с увеличилась в $2$ раза.
4) Путь, пройденный телом 1 за промежуток времени от $0$ до $18$~с, меньше пути, пройденного телом 2 за тот же промежуток времени.
5) Модуль импульса тела 2 в момент времени $t=10$~с равен $6$~\text{кг}$\cdot$\text{м/с}.

\kim{5}
\end{taskbasic}
\answer{14}



% === ЗАДАЧА 10 ===
% Тег: #МЕХАНИКА #КИМ_05 #ВАРИАНТ_09
\begin{taskbasic}[code={\label{task:newton2_v9_n5}}]
\textbf{Задача 10.} Мотоцикл с мотоциклистом проезжает верхнюю точку выпуклого моста, радиус кривизны которого равен $100$~м, двигаясь с постоянной по модулю скоростью, равной $54$~км/ч. Общая масса мотоцикла и мотоциклиста составляет $600$~кг. Из приведённого ниже списка выберите все верные утверждения, характеризующие движение мотоцикла с мотоциклистом.

1) Модуль центростремительного ускорения мотоцикла с мотоциклистом равен $2{,}25$~$м/с^2$.
2) Модуль силы тяжести, действующей на мотоцикл с мотоциклистом, равен $6$~кН.
3) Равнодействующая сил, действующих на мотоцикл с мотоциклистом, направлена вертикально вверх и перпендикулярна скорости мотоцикла.
4) Сила, с которой мотоцикл с мотоциклистом действует на мост, направлена вертикально вниз и равна по модулю $46{,}5$~кН.
5) Сила, с которой мост действует на мотоцикл с мотоциклистом, направлена вертикально вверх.

\kim{5}
\end{taskbasic}
\answer{25}

% === ЗАДАЧА 11 ===
% Тег: #МЕХАНИКА #КИМ_05 #ВАРИАНТ_10
\begin{taskbasic}[code={\label{task:newton2_v10_n5}}]
\textbf{Задача 11.} Мотоцикл с мотоциклистом проезжает нижнюю точку вогнутого моста, радиус кривизны которого равен $200$~м, двигаясь с постоянной по модулю скоростью, равной $72$~км/ч. Общая масса мотоцикла и мотоциклиста составляет $550$~кг. Из приведённого ниже списка выберите все верные утверждения, характеризующие движение мотоцикла с мотоциклистом.

1) Модуль равнодействующей сил, действующих на мотоцикл с мотоциклистом, равен $11000$~Н.
2) Равнодействующая сил, действующих на мотоцикл с мотоциклистом, направлена вертикально вниз.
3) Модуль силы тяжести, действующей на мотоциклиста с мотоциклом, на $1{,}1$~кН меньше модуля силы, с которой они действуют на мост.
4) Сила, с которой мотоцикл с мотоциклистом действует на мост, направлена вертикально вниз и равна по модулю $6600$~Н.
5) Сила, с которой мост действует на мотоцикл с мотоциклистом, сонаправлена с его центростремительным ускорением.

\kim{5}
\end{taskbasic}
\answer{35}

% ====================================================================
% ФАЙЛ: tasks/01_mechanics/newton_second_law_bank_part2.tex
% ТЕМА: ВТОРОЙ ЗАКОН НЬЮТОНА (ДИНАМИКА)
% ПАЧКА 2: ВАРИАНТЫ 11–20 (С ИСПРАВЛЕННЫМИ ЕДИНИЦАМИ И TIKZ-РИСУНКАМИ)
% ====================================================================

% === ЗАДАЧА 1 ===
% Тег: #МЕХАНИКА #КИМ_02 #ВАРИАНТ_11
\begin{taskbasic}[code={\label{task:newton2_v11_n2}}]
\textbf{Задача 1.} В инерциальной системе отсчёта некоторая сила сообщает телу массой $5$~кг ускорение, равное по модулю $4$~м/с$^{2}$. Телу какой массы вдвое большая сила сообщит в той же системе отсчёта ускорение $1$~м/с$^{2}$?

\kim{2}
\end{taskbasic}
\answer{40}

% === ЗАДАЧА 2 ===
% Тег: #МЕХАНИКА #КИМ_26 #ВАРИАНТ_11
\begin{taskbasic}[code={\label{task:newton2_v11_n26}}]
\textbf{Задача 2.} Если в установке первоначально покоящуюся тележку толкнуть влево, то она движется с ускорением $3$~м/с$^{2}$. Если же тележку толкнуть вправо, то она движется равномерно. Найдите массу $M$ тележки, если масса грузика на нити $m=150$~г. Массами блока и нити пренебречь. Нить нерастяжима. Модуль силы сопротивления движению тележки считать постоянным и одинаковым в обоих случаях, трением в оси блока пренебречь.
\nopagebreak\smallskip
\begin{center}
\begin{tikzpicture}[scale=1.2, every node/.style={font=\footnotesize}]
    % Стол (нижняя толстая плита и ножка)
    \draw[thick, fill=gray!20] (-2.5,-0.3) rectangle (3.0,0);
    \draw[thick, fill=gray!20] (1.8,-1.2) rectangle (2.0,-0.3);
    
    % Лабораторный трек / рельс (длинная серая плита на столе)
    \draw[thick, fill=gray!10] (-1.9,0) rectangle (2.9,0.15);
    % Левый упор (ограничитель) на треке
    \draw[thick, fill=gray!30] (-1.9,0.15) rectangle (-1.7,0.45);

    % Неподвижный блок на правом краю трека
    \draw[thick, fill=gray!40] (3.0,0.15) circle (0.12);
    \draw[thick, fill=gray!10] (3.0,0.15) circle (0.04);
    \fill[black] (3.0,0.15) circle (1pt);
    % Крепление блока к рельсу
    \draw[thick] (2.9,0.05) -- (3.0,0.15);

    % Тележка M (с ушками по краям)
    \draw[thick, fill=brandPrimary!10] (-0.8,0.22) rectangle (-0.2,0.42);
    \draw[thick, fill=brandPrimary!10] (-0.8,0.42) -- (-0.8,0.52) -- (-0.75,0.52) -- (-0.75,0.22);
    \draw[thick, fill=brandPrimary!10] (-0.2,0.42) -- (-0.2,0.52) -- (-0.25,0.52) -- (-0.25,0.22);
    \node[above] at (-0.5,0.52) {$M$};
    % Колёса тележки
    \draw[thick, fill=black!50] (-0.65,0.185) circle (0.04);
    \draw[thick, fill=black!50] (-0.35,0.185) circle (0.04);

    % Груз m (подвешенный вертикально)
    \draw[thick, fill=brandAccent!20] (2.85,-0.9) rectangle (3.15,-0.83);
    \node[right] at (3.2,-0.86) {$m$};

    % Горизонтальная нить (от правого края тележки к блоку)
    \draw[thick] (-0.2,0.42) -- (3.0,0.27);
    
    % Вертикальная нить (от блока к грузу m)
    \draw[thick] (3.12,0.15) -- (3.12,-0.83);
\end{tikzpicture}
\end{center}
\kim{26}
\end{taskbasic}
\answer{0{,}85~кг}

% === ЗАДАЧА 3 ===
% Тег: #МЕХАНИКА #КИМ_02 #ВАРИАНТ_12
\begin{taskbasic}[code={\label{task:newton2_v12_n2}}]
\textbf{Задача 3.} В инерциальной системе отсчёта некоторая сила сообщает телу массой $2{,}5$~кг ускорение, равное по модулю $6$~м/с$^{2}$. Определите модуль ускорения, которое сообщит та же сила в той же системе отсчёта телу массой $7{,}5$~кг.

\kim{2}
\end{taskbasic}
\answer{2}

% === ЗАДАЧА 4 ===
% Тег: #МЕХАНИКА #КИМ_26 #ВАРИАНТ_12
\begin{taskbasic}[code={\label{task:newton2_v12_n26}}]
\textbf{Задача 4.} Если в установке первоначально покоящуюся тележку толкнуть влево, то она движется с ускорением $2$~м/с$^{2}$. Если же тележку толкнуть вправо, то она движется равномерно. Найдите массу $m$ грузика на нити, если масса тележки $M=450$~г. Массами блока и нити пренебречь. Нить нерастяжима. Модуль силы сопротивления движению тележки считать постоянным и одинаковым в обоих случаях, трением в оси блока пренебречь.

\nopagebreak\smallskip
\begin{center}
\begin{tikzpicture}[scale=1.2, every node/.style={font=\footnotesize}]
    % Стол (нижняя толстая плита и ножка)
    \draw[thick, fill=gray!20] (-2.5,-0.3) rectangle (3.0,0);
    \draw[thick, fill=gray!20] (1.8,-1.2) rectangle (2.0,-0.3);
    
    % Лабораторный трек / рельс (длинная серая плита на столе)
    \draw[thick, fill=gray!10] (-1.9,0) rectangle (2.9,0.15);
    % Левый упор (ограничитель) на треке
    \draw[thick, fill=gray!30] (-1.9,0.15) rectangle (-1.7,0.45);

    % Неподвижный блок на правом краю трека
    \draw[thick, fill=gray!40] (3.0,0.15) circle (0.12);
    \draw[thick, fill=gray!10] (3.0,0.15) circle (0.04);
    \fill[black] (3.0,0.15) circle (1pt);
    % Крепление блока к рельсу
    \draw[thick] (2.9,0.05) -- (3.0,0.15);

    % Тележка M (с ушками по краям)
    \draw[thick, fill=brandPrimary!10] (-0.8,0.22) rectangle (-0.2,0.42);
    \draw[thick, fill=brandPrimary!10] (-0.8,0.42) -- (-0.8,0.52) -- (-0.75,0.52) -- (-0.75,0.22);
    \draw[thick, fill=brandPrimary!10] (-0.2,0.42) -- (-0.2,0.52) -- (-0.25,0.52) -- (-0.25,0.22);
    \node[above] at (-0.5,0.52) {$M$};
    % Колёса тележки
    \draw[thick, fill=black!50] (-0.65,0.185) circle (0.04);
    \draw[thick, fill=black!50] (-0.35,0.185) circle (0.04);

    % Груз m (подвешенный вертикально)
    \draw[thick, fill=brandAccent!20] (2.85,-0.9) rectangle (3.15,-0.83);
    \node[right] at (3.2,-0.86) {$m$};

    % Горизонтальная нить (от правого края тележки к блоку)
    \draw[thick] (-0.2,0.42) -- (3.0,0.27);
    
    % Вертикальная нить (от блока к грузу m)
    \draw[thick] (3.12,0.15) -- (3.12,-0.83);
\end{tikzpicture}
\end{center}

\kim{26}
\end{taskbasic}
\answer{0{,}045~кг}

% === ЗАДАЧА 5 ===
% Тег: #МЕХАНИКА #КИМ_06 #ВАРИАНТ_13
\begin{taskbasic}[code={\label{task:newton2_v13_n6}}]
\textbf{Задача 5.} Космический аппарат, обращающийся вокруг Венеры по круговой орбите, перешёл на другую круговую орбиту большего радиуса. Как изменились в результате этого перехода центростремительное ускорение, с которым аппарат движется по орбите, и его период обращения вокруг Венеры?

Для каждой величины определите соответствующий характер изменения:
1) увеличивается;
2) уменьшается;
3) не меняется.

\kim{6}
\end{taskbasic}
\answer{21}

% === ЗАДАЧА 6 ===
% Тег: #МЕХАНИКА #КИМ_26 #ВАРИАНТ_13
\begin{taskbasic}[code={\label{task:newton2_v13_n26}}]
\textbf{Задача 6.} На горизонтальном столе находится брусок массой $M=1$~кг, соединённый невесомой нерастяжимой нитью, перекинутой через гладкий невесомый блок, с грузом массой $m=500$~г. На брусок действует сила величиной $F=9$~Н, направленная под углом $\alpha=30^{\circ}$ к горизонту. В момент начала движения груз находится на расстоянии $L=40$~см от края стола. За какое время груз поднимется до края стола, если коэффициент трения между бруском и столом $\mu=0{,}3$? Трением в оси блока и трением о воздух пренебречь.

\nopagebreak\smallskip
\begin{center}
\begin{tikzpicture}[scale=1.2, every node/.style={font=\footnotesize}]
    % Поверхность стола и угловой обрыв (чистый fill=gray!5)
    \draw[thick, fill=gray!5] (-1.5,0) -- (3.5,0) -- (3.5,-0.1) -- (-1.4,-0.1) -- (-1.4,-2.2) -- (-1.5,-2.2) -- cycle;
    % Штриховка обрыва стола
    \foreach \y in {-0.2,-0.4,-0.6,-0.8,-1.0,-1.2,-1.4,-1.6,-1.8,-2.0}
        \draw[thin] (-1.5,\y) -- (-1.4,\y+0.1);

    % Неподвижный блок на углу стола
    \draw[thick, fill=gray!30] (-1.5,0) circle (0.18);
    \draw[thick, fill=gray!10] (-1.5,0) circle (0.06);
    \fill[black] (-1.5,0) circle (1.2pt);

    % Брусок M (на столе)
    \draw[thick, fill=brandPrimary!10] (0.5,0) rectangle (1.7,0.7);
    \node at (1.1,0.35) {$M$};
    
    % Сила F, приложенная к бруску M
    \draw[-{Stealth[scale=0.8]}, ultra thick, brandAccent] (1.7,0.35) -- (2.7,0.85) node[above] {$\vec{F}$};
    \draw[dashed, thick] (1.7,0.35) -- (2.7,0.35);
    \draw (2.1,0.35) arc (0:26.5:0.4);
    \node at (2.25,0.47) {$\alpha$};
    
    % Груз m (висит вертикально)
    \draw[thick, fill=brandAccent!20] (-1.9,-1.5) rectangle (-1.45,-1.0);
    \node at (-1.675,-1.25) {$m$};
    
    % Горизонтальный участок нити (от блока к бруску M)
    \draw[thick] (-1.5,0.18) -- (0.5,0.18);
    
    % Вертикальный участок нити (от блока к грузу m)
    \draw[thick] (-1.68,0) -- (-1.68,-1.0);
    
    % Выносные линии и стрелка для расстояния L
    \draw[gray] (-1.4,0) -- (-1.0,0);
    \draw[gray] (-1.45,-1.0) -- (-1.0,-1.0);
    \draw[<->, gray] (-1.1,0) -- (-1.1,-1.0) node[midway, right, black] {$L$};
\end{tikzpicture}
\end{center}

\kim{26}
\end{taskbasic}
\answer{1~с}

% === ЗАДАЧА 7 ===
% Тег: #МЕХАНИКА #КИМ_06 #ВАРИАНТ_14
\begin{taskbasic}[code={\label{task:newton2_v14_n6}}]
\textbf{Задача 7.} Космический аппарат, обращающийся вокруг Луны по круговой орбите, перешёл на другую круговую орбиту меньшего радиуса. Как изменились в результате этого перехода кинетическая энергия аппарата и его частота обращения вокруг Луны?

Для каждой величины определите соответствующий характер изменения:
1) увеличилась;
2) уменьшилась;
3) не изменилась.

\kim{6}
\end{taskbasic}
\answer{11}

% === ЗАДАЧА 8 ===
% Тег: #МЕХАНИКА #КИМ_26 #ВАРИАНТ_14
\begin{taskbasic}[code={\label{task:newton2_v14_n26}}]
\textbf{Задача 8.} На горизонтальном столе находится брусок массой $M=1$~кг, соединённый невесомой нерастяжимой нитью, перекинутой через гладкий невесомый блок, с грузом массой $m=500$~г. На брусок действует сила величиной $F=8$~Н, направленная под углом $\alpha=30^{\circ}$ к горизонту. В момент начала движения груз находится на расстоянии $L=1$~м от края стола. Каков коэффициент трения между бруском и столом, если груз поднимается до края стола за время $t=2$~с? Трением в оси блока и трением о воздух пренебречь.

\nopagebreak\smallskip
\begin{center}
\begin{tikzpicture}[scale=1.2, every node/.style={font=\footnotesize}]
    % Поверхность стола и угловой обрыв (чистый fill=gray!5)
    \draw[thick, fill=gray!5] (-1.5,0) -- (3.5,0) -- (3.5,-0.1) -- (-1.4,-0.1) -- (-1.4,-2.2) -- (-1.5,-2.2) -- cycle;
    % Штриховка обрыва стола
    \foreach \y in {-0.2,-0.4,-0.6,-0.8,-1.0,-1.2,-1.4,-1.6,-1.8,-2.0}
        \draw[thin] (-1.5,\y) -- (-1.4,\y+0.1);

    % Неподвижный блок на углу стола
    \draw[thick, fill=gray!30] (-1.5,0) circle (0.18);
    \draw[thick, fill=gray!10] (-1.5,0) circle (0.06);
    \fill[black] (-1.5,0) circle (1.2pt);

    % Брусок M (на столе)
    \draw[thick, fill=brandPrimary!10] (0.5,0) rectangle (1.7,0.7);
    \node at (1.1,0.35) {$M$};
    
    % Сила F, приложенная к бруску M
    \draw[-{Stealth[scale=0.8]}, ultra thick, brandAccent] (1.7,0.35) -- (2.7,0.85) node[above] {$\vec{F}$};
    \draw[dashed, thick] (1.7,0.35) -- (2.7,0.35);
    \draw (2.1,0.35) arc (0:26.5:0.4);
    \node at (2.25,0.47) {$\alpha$};
    
    % Груз m (висит вертикально)
    \draw[thick, fill=brandAccent!20] (-1.9,-1.5) rectangle (-1.45,-1.0);
    \node at (-1.675,-1.25) {$m$};
    
    % Горизонтальный участок нити (от блока к бруску M)
    \draw[thick] (-1.5,0.18) -- (0.5,0.18);
    
    % Вертикальный участок нити (от блока к грузу m)
    \draw[thick] (-1.68,0) -- (-1.68,-1.0);
    
    % Выносные линии и стрелка для расстояния L
    \draw[gray] (-1.4,0) -- (-1.0,0);
    \draw[gray] (-1.45,-1.0) -- (-1.0,-1.0);
    \draw[<->, gray] (-1.1,0) -- (-1.1,-1.0) node[midway, right, black] {$L$};
\end{tikzpicture}
\end{center}

\kim{26}
\end{taskbasic}
\answer{0{,}2}

% === ЗАДАЧА 9 ===
% Тег: #МЕХАНИКА #КИМ_06 #ВАРИАНТ_15
\begin{taskbasic}[code={\label{task:newton2_v15_n6}}]
\textbf{Задача 9.} На рисунке показан график зависимости координаты $x$ тела, движущегося вдоль оси $Ox$, от времени $t$ (парабола). Графики А и Б представляют собой зависимости физических величин, характеризующих движение этого тела, от времени $t$.
\begin{center}
\begin{minipage}[b]{0.45\linewidth}
\centering
\begin{tikzpicture}[scale=0.7, every node/.style={font=\footnotesize}]
    \node[above] at (1.2,2.0) {А)};
    \draw[-{Stealth[scale=0.8]}, black, thick] (0,-1.2) -- (0,1.5) node[above] {};
    \draw[-{Stealth[scale=0.8]}, black, thick] (0,0) -- (3,0) node[right] {$t$};
    \draw[ultra thick, brandPrimary] (0,-0.8) -- (2,0) -- (3,0.4);
    \node[below] at (2,0) {$t_1$};
    \node[below left] at (0,0) {0};
\end{tikzpicture}
\end{minipage}
\begin{minipage}[b]{0.45\linewidth}
\centering
\begin{tikzpicture}[scale=0.7, every node/.style={font=\footnotesize}]
    \node[above] at (1.2,2.0) {Б)};
    \draw[-{Stealth[scale=0.8]}, black, thick] (0,0) -- (0,2.2) node[above] {};
    \draw[-{Stealth[scale=0.8]}, black, thick] (0,0) -- (3,0) node[right] {$t$};
    \draw[ultra thick, brandPrimary] (0,1.0) -- (3,1.0);
    \node[below left] at (0,0) {0};
\end{tikzpicture}
\end{minipage}
\end{center}
Установите соответствие между графиками и физическими величинами, зависимость которых от времени эти графики могут представлять.
1) проекция скорости тела на ось $Ox$
2) проекция перемещения тела на ось $Ox$
3) кинетическая энергия тела
4) модуль равнодействующей сил, действующих на тело

\kim{6}
\end{taskbasic}
\answer{14}

% === ЗАДАЧА 10 ===
% Тег: #МЕХАНИКА #КИМ_06 #ВАРИАНТ_16
\begin{taskbasic}[code={\label{task:newton2_v16_n6}}]
\textbf{Задача 10.} На рисунке показан график зависимости координаты $x$ тела, движущегося вдоль оси $Ox$, от времени $t$ (парабола). Графики А и Б представляют собой зависимости физических величин, характеризующих движение этого тела, от времени $t$.
\begin{center}
\begin{minipage}[b]{0.45\linewidth}
\centering
\begin{tikzpicture}[scale=0.7, every node/.style={font=\footnotesize}]
    \node[above] at (1.2,2.0) {А)};
    \draw[-{Stealth[scale=0.8]}, black, thick] (0,0) -- (0,2.2) node[above] {};
    \draw[-{Stealth[scale=0.8]}, black, thick] (0,0) -- (3,0) node[right] {$t$};
    \draw[ultra thick, brandPrimary, domain=0:2.6, samples=100] plot (\x, {1.2 - 1.2*\x + 0.4*\x*\x});
    \node[below] at (1.5,0) {$t_1$};
    \node[below left] at (0,0) {0};
\end{tikzpicture}
\end{minipage}
\begin{minipage}[b]{0.45\linewidth}
\centering
\begin{tikzpicture}[scale=0.7, every node/.style={font=\footnotesize}]
    \node[above] at (1.2,2.0) {Б)};
    \draw[-{Stealth[scale=0.8]}, black, thick] (0,0) -- (0,2.2) node[above] {};
    \draw[-{Stealth[scale=0.8]}, black, thick] (0,0) -- (3,0) node[right] {$t$};
    \draw[ultra thick, brandPrimary] (0,1.0) -- (3,1.0);
    \node[below left] at (0,0) {0};
\end{tikzpicture}
\end{minipage}
\end{center}
Установите соответствие между графиками и физическими величинами, зависимость которых от времени эти графики могут представлять.
1) проекция на ось $Ox$ скорости тела
2) проекция перемещения на ось $Ox$ тела
3) кинетическая энергия тела
4) модуль равнодействующей сил, действующих на тело

\kim{6}
\end{taskbasic}
\answer{34}

% === ЗАДАЧА 11 ===
% Тег: #МЕХАНИКА #КИМ_26 #ВАРИАНТ_17
\begin{taskbasic}[code={\label{task:newton2_v17_n26}}]
\textbf{Задача 11.} На гладком горизонтальном столе лежит доска массой $M=1$~кг и длиной $L=50$~см. На левом краю доски находится маленький брусок массой $m=200$~г. Брусок и доска связаны невесомой нерастяжимой нитью, перекинутой через невесомый гладкий блок, закреплённый на стене (отрезки нити, не лежащие на блоке, горизонтальны и параллельны). Коэффициент трения между бруском и доской $\mu=0{,}3$. Брусок начинают тянуть вправо постоянной силой $\vec{F}$, параллельной горизонтальным отрезкам нити. Через время $t=2$~с после начала движения брусок соскальзывает с доски. Определите модуль силы $\vec{F}$. Размерами бруска пренебречь.

\nopagebreak\smallskip
\begin{center}
\begin{tikzpicture}[scale=1.2, every node/.style={font=\footnotesize}]
    % Поверхность пола (горизонтальная линия стола)
    \draw[thick, fill=gray!5] (-2.2,0) -- (4.2,0) -- (4.2,-0.1) -- (-2.2,-0.1) -- cycle;
    
    % Левая вертикальная стена (опора для блока)
    \draw[thick, fill=gray!20] (-2.2,0) rectangle (-2.0,2.0);
    % Штриховка стены
    \foreach \y in {0.1,0.3,0.5,0.7,0.9,1.1,1.3,1.5,1.7,1.9}
        \draw[thin] (-2.2,\y) -- (-2.0,\y-0.1);

    % Доска M (большая нижняя доска)
    \draw[thick, fill=brandPrimary!10] (1.0,0) rectangle (3.4,0.5);
    \node at (2.2,0.25) {$M$};
    
    % Брусок m (маленький верхний брусок, на левом краю доски M)
    \draw[thick, fill=brandAccent!20] (1.0,0.5) rectangle (1.7,1.0);
    \node at (1.35,0.75) {$m$};
    
    % Неподвижный блок на левой стене
    \draw[thick, fill=gray!30] (-1.4,0.65) circle (0.25);
    \draw[thick, fill=gray!10] (-1.4,0.65) circle (0.08);
    \fill[black] (-1.4,0.65) circle (1.5pt);
    
    % Крепление блока к левой стене
    \draw[thick] (-2.0,0.65) -- (-1.4,0.65);
    
    % Верхняя нить (от блока к бруску m)
    \draw[thick] (-1.4,0.9) -- (1.0,0.9);
    
    % Нижняя нить (от блока к доске M)
    \draw[thick] (-1.4,0.4) -- (1.0,0.4);
    
    % Сила F, тянущая брусок m вправо
    \draw[-{Stealth[scale=0.8]}, ultra thick, brandAccent] (1.7,0.75) -- (3.0,0.75) node[above] {$\vec{F}$};
    
    % Выносные линии и стрелка для длины L доски M
    \draw[gray] (1.0,0) -- (1.0,-0.4);
    \draw[gray] (3.4,0) -- (3.4,-0.4);
    \draw[<->, gray] (1.0,-0.3) -- (3.4,-0.3) node[midway, below, black] {$L$};
\end{tikzpicture}
\end{center}

\kim{26}
\end{taskbasic}
\answer{1{,}35~Н}

% === ЗАДАЧА 12 ===
% Тег: #МЕХАНИКА #КИМ_26 #ВАРИАНТ_18
\begin{taskbasic}[code={\label{task:newton2_v18_n26}}]
\textbf{Задача 12.} На гладком горизонтальном столе лежит доска массой $M=1$~кг и длиной $L=50$~см. На левом краю доски находится маленький брусок массой $m=200$~г. Брусок и доска связаны невесомой нерастяжимой нитью, перекинутой через невесомый гладкий блок, закреплённый на стене (отрезки нити, не лежащие на блоке, горизонтальны и параллельны). Коэффициент трения между бруском и доской $\mu=0{,}2$. Брусок начинают тянуть вправо с постоянной силой $\vec{F}$, параллельной горизонтальным отрезкам нити. Модуль этой силы $F=1{,}2$~Н. Через какое время после начала движения брусок соскользнёт с доски? Размерами бруска пренебречь.

\nopagebreak\smallskip
\begin{center}
\begin{tikzpicture}[scale=1.2, every node/.style={font=\footnotesize}]
    % Поверхность пола (горизонтальная линия стола)
    \draw[thick, fill=gray!5] (-2.2,0) -- (4.2,0) -- (4.2,-0.1) -- (-2.2,-0.1) -- cycle;
    
    % Левая вертикальная стена (опора для блока)
    \draw[thick, fill=gray!20] (-2.2,0) rectangle (-2.0,2.0);
    % Штриховка стены
    \foreach \y in {0.1,0.3,0.5,0.7,0.9,1.1,1.3,1.5,1.7,1.9}
        \draw[thin] (-2.2,\y) -- (-2.0,\y-0.1);

    % Доска M (большая нижняя доска)
    \draw[thick, fill=brandPrimary!10] (1.0,0) rectangle (3.4,0.5);
    \node at (2.2,0.25) {$M$};
    
    % Брусок m (маленький верхний брусок, на левом краю доски M)
    \draw[thick, fill=brandAccent!20] (1.0,0.5) rectangle (1.7,1.0);
    \node at (1.35,0.75) {$m$};
    
    % Неподвижный блок на левой стене
    \draw[thick, fill=gray!30] (-1.4,0.65) circle (0.25);
    \draw[thick, fill=gray!10] (-1.4,0.65) circle (0.08);
    \fill[black] (-1.4,0.65) circle (1.5pt);
    
    % Крепление блока к левой стене
    \draw[thick] (-2.0,0.65) -- (-1.4,0.65);
    
    % Верхняя нить (от блока к бруску m)
    \draw[thick] (-1.4,0.9) -- (1.0,0.9);
    
    % Нижняя нить (от блока к доске M)
    \draw[thick] (-1.4,0.4) -- (1.0,0.4);
    
    % Сила F, тянущая брусок m вправо
    \draw[-{Stealth[scale=0.8]}, ultra thick, brandAccent] (1.7,0.75) -- (3.0,0.75) node[above] {$\vec{F}$};
    
    % Выносные линии и стрелка для длины L доски M
    \draw[gray] (1.0,0) -- (1.0,-0.4);
    \draw[gray] (3.4,0) -- (3.4,-0.4);
    \draw[<->, gray] (1.0,-0.3) -- (3.4,-0.3) node[midway, below, black] {$L$};
\end{tikzpicture}
\end{center}

\kim{26}
\end{taskbasic}
\answer{1{,}22~с}

% === ЗАДАЧА 13 ===
% Тег: #МЕХАНИКА #КИМ_02 #ВАРИАНТ_19
\begin{taskbasic}[code={\label{task:newton2_v19_n2}}]
\textbf{Задача 13.} В инерциальной системе отсчёта сила величиной $35$~Н сообщает телу массой $5$~кг некоторое ускорение. Сила какой величины сообщит телу массой $9$~кг в этой же системе отсчёта такое же ускорение?

\kim{2}
\end{taskbasic}
\answer{63}

% === ЗАДАЧА 14 ===
% Тег: #МЕХАНИКА #КИМ_26 #ВАРИАНТ_19
\begin{taskbasic}[code={\label{task:newton2_v19_n26}}]
\textbf{Задача 14.} На горизонтальном неподвижном столе лежит доска, на которой находится маленький брусок массой $m=400$~г. Брусок и доска связаны невесомой нерастяжимой нитью, перекинутой через невесомый блок, закреплённый на стене (отрезки нити, не лежащие на блоке, горизонтальны и параллельны). Коэффициент трения между бруском и доской $\mu_1=0{,}5$, между столом и доской $\mu_2=0{,}3$. Доску тянут вправо постоянной горизонтальной силой $F$, параллельной горизонтальным отрезкам нити, $F=12$~Н. Чему равна масса доски $M$, если модуль ускорения бруска относительно стола $a=2$~м/с$^{2}$? Трением в оси блока пренебречь.

\nopagebreak\smallskip
\begin{center}
\begin{tikzpicture}[scale=1.2, every node/.style={font=\footnotesize}]
    % Поверхность пола (горизонтальная линия стола)
    \draw[thick, fill=gray!5] (-2.2,0) -- (4.2,0) -- (4.2,-0.1) -- (-2.2,-0.1) -- cycle;
    
    % Левая вертикальная стена (опора для блока)
    \draw[thick, fill=gray!20] (-2.2,0) rectangle (-2.0,2.0);
    % Штриховка стены
    \foreach \y in {0.1,0.3,0.5,0.7,0.9,1.1,1.3,1.5,1.7,1.9}
        \draw[thin] (-2.2,\y) -- (-2.0,\y-0.1);

    % Доска M (большая нижняя доска)
    \draw[thick, fill=brandPrimary!10] (0.2,0) rectangle (2.6,0.5);
    \node at (1.4,0.25) {$M$};
    
    % Брусок m (маленький верхний брусок)
    \draw[thick, fill=brandAccent!20] (1.9,0.5) rectangle (2.6,1.0);
    \node at (2.25,0.75) {$m$};
    
    % Неподвижный блок на левой стене
    \draw[thick, fill=gray!30] (-1.4,0.65) circle (0.25);
    \draw[thick, fill=gray!10] (-1.4,0.65) circle (0.08);
    \fill[black] (-1.4,0.65) circle (1.5pt);
    
    % Крепление блока к левой стене
    \draw[thick] (-2.0,0.65) -- (-1.4,0.65);
    
    % Верхняя нить (от блока к бруску m)
    \draw[thick] (-1.4,0.9) -- (1.9,0.9);
    
    % Нижняя нить (от блока к доске M)
    \draw[thick] (-1.4,0.4) -- (0.2,0.4);
    
    % Сила F, тянущая доску M вправо
    \draw[-{Stealth[scale=0.8]}, ultra thick, brandPrimary] (2.6,0.25) -- (3.6,0.25) node[above] {$\vec{F}$};
\end{tikzpicture}
\end{center}

\kim{26}
\end{taskbasic}
\answer{1{,}2~кг}

% === ЗАДАЧА 15 ===
% Тег: #МЕХАНИКА #КИМ_02 #ВАРИАНТ_20
\begin{taskbasic}[code={\label{task:newton2_v20_n2}}]
\textbf{Задача 15.} В инерциальной системе отсчёта сила величиной $70$~Н сообщает телу массой $10$~кг некоторое ускорение. Телу какой массы сила величиной $56$~Н сообщит в этой же системе отсчёта вдвое большее ускорение?

\kim{2}
\end{taskbasic}
\answer{4}

% === ЗАДАЧА 16 ===
% Тег: #МЕХАНИКА #КИМ_26 #ВАРИАНТ_20
\begin{taskbasic}[code={\label{task:newton2_v20_n26}}]
\textbf{Задача 16.} На горизонтальном неподвижном столе лежит доска массой $M=0{,}8$~кг. На доске находится маленький брусок массой $m=200$~г. Брусок и доска связаны невесомой нерастяжимой нитью, перекинутой через невесомый блок, закреплённый на стене (отрезки нити, не лежащие на блоке, горизонтальны и параллельны). Коэффициент трения между бруском и доской $\mu_1=0{,}5$, между столом и доской $\mu_2=0{,}3$. Доску тянут вправо постоянной силой $\vec{F}$, параллельной горизонтальным отрезкам нити. Чему равен модуль силы $\vec{F}$, если модуль ускорения бруска относительно стола $a=1$~м/с$^{2}$? Трением в оси блока пренебречь.

\nopagebreak\smallskip
\begin{center}
\begin{tikzpicture}[scale=1.2, every node/.style={font=\footnotesize}]
    % Поверхность пола (горизонтальная линия стола)
    \draw[thick, fill=gray!5] (-2.2,0) -- (4.2,0) -- (4.2,-0.1) -- (-2.2,-0.1) -- cycle;
    
    % Левая вертикальная стена (опора для блока)
    \draw[thick, fill=gray!20] (-2.2,0) rectangle (-2.0,2.0);
    % Штриховка стены
    \foreach \y in {0.1,0.3,0.5,0.7,0.9,1.1,1.3,1.5,1.7,1.9}
        \draw[thin] (-2.2,\y) -- (-2.0,\y-0.1);

    % Доска M (большая нижняя доска)
    \draw[thick, fill=brandPrimary!10] (0.2,0) rectangle (2.6,0.5);
    \node at (1.4,0.25) {$M$};
    
    % Брусок m (маленький верхний брусок)
    \draw[thick, fill=brandAccent!20] (1.9,0.5) rectangle (2.6,1.0);
    \node at (2.25,0.75) {$m$};
    
    % Неподвижный блок на левой стене
    \draw[thick, fill=gray!30] (-1.4,0.65) circle (0.25);
    \draw[thick, fill=gray!10] (-1.4,0.65) circle (0.08);
    \fill[black] (-1.4,0.65) circle (1.5pt);
    
    % Крепление блока к левой стене
    \draw[thick] (-2.0,0.65) -- (-1.4,0.65);
    
    % Верхняя нить (от блока к бруску m)
    \draw[thick] (-1.4,0.9) -- (1.9,0.9);
    
    % Нижняя нить (от блока к доске M)
    \draw[thick] (-1.4,0.4) -- (0.2,0.4);
    
    % Сила F, тянущая доску M вправо
    \draw[-{Stealth[scale=0.8]}, ultra thick, brandPrimary] (2.6,0.25) -- (3.6,0.25) node[above] {$\vec{F}$};
\end{tikzpicture}
\end{center}

\kim{26}
\end{taskbasic}
\answer{6~Н}

% ====================================================================
% ФАЙЛ: tasks/01_mechanics/newton_second_law_bank_part3.tex
% ТЕМА: ВТОРОЙ ЗАКОН НЬЮТОНА (ДИНАМИКА И СВЯЗАННЫЕ ТЕЛА)
% ПАЧКА 3: ВАРИАНТЫ 21–30
% ====================================================================

% === ЗАДАЧА 1 ===
% Тег: #МЕХАНИКА #КИМ_02 #ВАРИАНТ_22
\begin{taskbasic}[code={\label{task:newton2_v22_n2}}]
\textbf{Задача 1.} В инерциальной системе отсчёта сила $F_1 = 4$~Н сообщает телу массой $m$ ускорение $a$. Каков модуль силы $F_2$, которая сообщает телу массой $3m$ в этой системе отсчёта такое же ускорение?

\kim{2}
\end{taskbasic}
\answer{12~Н}

% === ЗАДАЧА 2 ===
% Тег: #МЕХАНИКА #КИМ_26 #ВАРИАНТ_22
\begin{taskbasic}[code={\label{task:newton2_v22_n26}}]
\textbf{Задача 2.} К бруску массой $M=2$~кг прикреплён лёгкий блок, через который переброшена лёгкая нерастяжимая нить. Один конец нити привязан к стене, а к другому прикреплено тело массой $m=0{,}75$~кг. На брусок действует сила $F=10$~Н. Определите ускорение бруска. Свободные куски нити горизонтальны и лежат в одной вертикальной плоскости, тела двигаются вдоль одной прямой. Массой блока и нити, а также трением пренебречь.

\nopagebreak\smallskip
\begin{center}
\begin{tikzpicture}[scale=1.2, every node/.style={font=\footnotesize}]
    % Поверхность стола (земля с заливкой)
    \draw[thick, fill=gray!5] (-1.0,0) -- (4.2,0) -- (4.2,-0.1) -- (-1.0,-0.1) -- cycle;
    
    % Правая вертикальная стена (опора для нити)
    \draw[thick, fill=gray!20] (3.6,0) rectangle (3.7,1.3);
    % Штриховка стены
    \foreach \y in {0.1,0.3,0.5,0.7,0.9,1.1}
        \draw[thin] (3.6,\y) -- (3.7,\y+0.1);

    % Брусок M (большой левый брусок)
    \draw[thick, fill=brandPrimary!10] (0,0) rectangle (0.8,1.2);
    \node at (0.4,0.6) {$M$};
    
    % Крепление подвижного блока к бруску M
    \draw[thick, fill=gray!30] (0.8,0.6) rectangle (1.2,0.7);
    
    % Подвижный блок
    \draw[thick, fill=gray!10] (1.4,0.65) circle (0.25);
    \fill[black] (1.4,0.65) circle (1.5pt);
    
    % Груз m (маленький правый брусок)
    \draw[thick, fill=brandAccent!20] (2.4,0) rectangle (3.1,0.5);
    \node at (2.75,0.25) {$m$};
    
    % Верхняя нить (от блока к стене)
    \draw[thick] (1.4,0.9) -- (3.6,0.9);
    
    % Нижняя нить (от блока к грузу m)
    \draw[thick] (1.4,0.4) -- (2.4,0.4);
    
    % Сила F, тянущая брусок M влево
    \draw[-{Stealth[scale=0.8]}, ultra thick, brandPrimary] (0,0.6) -- (-0.8,0.6) node[above] {$\vec{F}$};
\end{tikzpicture}
\end{center}

\kim{26}
\end{taskbasic}
\answer{2~м/с$^{2}$}

% === ЗАДАЧА 3 ===
% Тег: #МЕХАНИКА #КИМ_02 #ВАРИАНТ_24
\begin{taskbasic}[code={\label{task:newton2_v24_n2}}]
\textbf{Задача 3.} Модуль сил гравитационного притяжения между двумя шарами, находящимися на расстоянии $2$~м друг от друга, равен $9$~нН. Каков будет модуль сил притяжения между ними, если расстояние увеличить до $6$~м?

\kim{2}
\end{taskbasic}
\answer{1~нН}

% === ЗАДАЧА 4 ===
% Тег: #МЕХАНИКА #КИМ_02 #ВАРИАНТ_25
\begin{taskbasic}[code={\label{task:newton2_v25_n2}}]
\textbf{Задача 4.} На графике приведена зависимость ускорения $a$ бруска, скользящего без трения по горизонтальной поверхности, от величины приложенной к нему горизонтальной силы $\vec{F}$. Систему отсчёта считать инерциальной. Чему равна масса бруска?
\begin{center}
\begin{tikzpicture}[x=0.6cm, y=10cm, every node/.style={font=\footnotesize}]
    \draw[xstep=1, ystep=0.05, gray!30, thin] (0,0) grid (6,0.25);
    \draw[-{Stealth[scale=0.8]}, black, thick] (0,0) -- (6.5,0) node[right] {$F, \text{ Н}$};
    \draw[-{Stealth[scale=0.8]}, black, thick] (0,0) -- (0,0.27) node[above] {$a, \text{ м/с}^{2}$};
    \foreach \x in {2,4,6} \node[below] at (\x,0) {\x};
    \node[left] at (0,0.1) {$0{,}1$};
    \node[left] at (0,0.2) {$0{,}2$};
    \node[below left] at (0,0) {0};
    \draw[ultra thick, brandPrimary] (0,0) -- (5,0.2);
    \fill[black] (5,0.2) circle (1.5pt);
\end{tikzpicture}
\end{center}

\kim{2}
\end{taskbasic}
\answer{25~кг}

% === ЗАДАЧА 5 ===
% Тег: #МЕХАНИКА #КИМ_02 #ВАРИАНТ_26
\begin{taskbasic}[code={\label{task:newton2_v26_n2}}]
\textbf{Задача 5.} На графике приведена зависимость модуля горизонтальной силы $\vec{F}$, приложенной к бруску, скользящему без трения по горизонтальной поверхности, от модуля его ускорения $a$. Систему отсчёта считать инерциальной. Чему равна масса бруска?
\begin{center}
\begin{tikzpicture}[x=6cm, y=0.25cm, every node/.style={font=\footnotesize}]
    \draw[xstep=0.1, ystep=2, gray!30, thin] (0,0) grid (0.6,12);
    \draw[-{Stealth[scale=0.8]}, black, thick] (0,0) -- (0.68,0) node[right] {$a, \text{ м/с}^{2}$};
    \draw[-{Stealth[scale=0.8]}, black, thick] (0,0) -- (0,13) node[above] {$F, \text{ Н}$};
    \foreach \y in {3,6,12} \node[left] at (0,\y) {\y};
    \node[below] at (0.2,0) {$0{,}2$};
    \node[below] at (0.4,0) {$0{,}4$};
    \node[below] at (0.6,0) {$0{,}6$};
    \node[below left] at (0,0) {0};
    \draw[ultra thick, brandPrimary] (0,0) -- (0.6,12);
    \fill[black] (0.6,12) circle (1.5pt);
\end{tikzpicture}
\end{center}

\kim{2}
\end{taskbasic}
\answer{20~кг}

% === ЗАДАЧА 6 ===
% Тег: #МЕХАНИКА #КИМ_22 #ВАРИАНТ_27
\begin{taskbasic}[code={\label{task:newton2_v27_n22}}]
\textbf{Задача 6.} Груз массой $200$~г подвешен на пружине жёсткостью $100$~Н/м к потолку лифта. Лифт равноускоренно движется вниз, набирая скорость. Каково ускорение лифта, если удлинение пружины постоянно и равно $1{,}5$~см?

\kim{22}
\end{taskbasic}
\answer{2,5~м/с$^{2}$}

% === ЗАДАЧА 7 ===
% Тег: #МЕХАНИКА #КИМ_22 #ВАРИАНТ_28
\begin{taskbasic}[code={\label{task:newton2_v28_n22}}]
\textbf{Задача 7.} Груз массой $200$~г подвешен на пружине к потолку лифта. Лифт равноускоренно движется вниз с ускорением $2$~м/с$^{2}$, набирая скорость. Какова жёсткость пружины, если удлинение её постоянно и равно $1$~см?

\kim{22}
\end{taskbasic}
\answer{160~Н/м}

% === ЗАДАЧА 8 ===
% Тег: #МЕХАНИКА #КИМ_22 #ВАРИАНТ_29
\begin{taskbasic}[code={\label{task:newton2_v29_n22}}]
\textbf{Задача 8.} Два груза, связанных нерастяжимой и невесомой нитью, движутся по гладкой горизонтальной поверхности под действием горизонтальной силы $\vec{F}$, приложенной к грузу массой $M_1$. Максимальная сила $F$, при которой нить ещё не обрывается, равна $18$~Н. Известно, что нить может выдержать нагрузку не более $10$~Н. Чему равна масса $M_1$ первого груза, если масса второго равна $M_2 = 3$~кг?

\nopagebreak\smallskip
\begin{center}
\begin{tikzpicture}[scale=1.2, every node/.style={font=\footnotesize}]
    % Поверхность стола (чистый fill=gray!5)
    \draw[thick, fill=gray!5] (-2.2,0) -- (3.2,0) -- (3.2,-0.1) -- (-2.2,-0.1) -- cycle;
    
    % Левый груз m_2 (он же M_2 в тексте)
    \draw[thick, fill=brandAccent!20] (-1.3,0) rectangle (-0.3,0.5);
    \node at (-0.8,0.25) {$m_2$};
    
    % Правый груз m_1 (он же M_1 в тексте)
    \draw[thick, fill=brandPrimary!10] (0.8,0) rectangle (1.8,0.5);
    \node at (1.3,0.25) {$m_1$};
    
    % Нить между грузами
    \draw[thick] (-0.3,0.25) -- (0.8,0.25);
    
    % Сила F, приложенная к первому грузу
    \draw[-{Stealth[scale=0.8]}, ultra thick, brandPrimary] (1.8,0.25) -- (2.8,0.25) node[above] {$\vec{F}$};
\end{tikzpicture}
\end{center}
\kim{22}
\end{taskbasic}
\answer{2,4~кг}



% === ЗАДАЧА 10 ===
% Тег: #МЕХАНИКА #КИМ_02 #ВАРИАНТ_30
\begin{taskbasic}[code={\label{task:newton2_v30_n2}}]
\textbf{Задача 10.} В инерциальной системе отсчёта сила, модуль которой равен $16$~Н, сообщает телу массой $m$ ускорение $a$. Каков модуль силы, которая сообщает телу массой $4m$ в этой системе отсчёта ускорение $\frac{a}{2}$?

\kim{2}
\end{taskbasic}
\answer{32~Н}

% === ЗАДАЧА 11 ===
% Тег: #МЕХАНИКА #КИМ_22 #ВАРИАНТ_30
\begin{taskbasic}[code={\label{task:newton2_v30_n22}}]
\textbf{Задача 11.} Два груза, связанных нерастяжимой и невесомой нитью, движутся по гладкой горизонтальной поверхности под действием горизонтальной силы $\vec{F}$, приложенной к грузу массой $M_1 = 2$~кг. Максимальная сила $F$, при которой нить ещё не обрывается, равна $18$~Н. Известно, что нить может выдержать нагрузку не более $10$~Н. Чему равна масса $M_2$ второго груза?
\nopagebreak\smallskip
\begin{center}
\begin{tikzpicture}[scale=1.2, every node/.style={font=\footnotesize}]
    % Поверхность стола (чистый fill=gray!5)
    \draw[thick, fill=gray!5] (-2.2,0) -- (3.2,0) -- (3.2,-0.1) -- (-2.2,-0.1) -- cycle;
    
    % Левый груз m_2 (он же M_2 в тексте)
    \draw[thick, fill=brandAccent!20] (-1.3,0) rectangle (-0.3,0.5);
    \node at (-0.8,0.25) {$m_2$};
    
    % Правый груз m_1 (он же M_1 в тексте)
    \draw[thick, fill=brandPrimary!10] (0.8,0) rectangle (1.8,0.5);
    \node at (1.3,0.25) {$m_1$};
    
    % Нить между грузами
    \draw[thick] (-0.3,0.25) -- (0.8,0.25);
    
    % Сила F, приложенная к первому грузу
    \draw[-{Stealth[scale=0.8]}, ultra thick, brandPrimary] (1.8,0.25) -- (2.8,0.25) node[above] {$\vec{F}$};
\end{tikzpicture}
\end{center}

\kim{22}
\end{taskbasic}
\answer{2,5~кг}
% === ЗАДАЧА X ===
% Тег: #МЕХАНИКА #КИМ_26 #ДИНАМИКА_НАКЛОННОЙ_ПЛОСКОСТИ
\begin{taskbasic}[code={\label{task:elastic_inclined_connected_n26}}]
\textbf{Задача X.} По неподвижной гладкой наклонной плоскости с углом $\alpha = 30^{\circ}$ движутся два одинаковых бруска массой $m = 0{,}25$~кг каждый, скреплённые между собой лёгкой пружиной с жёсткостью $k = 100$~Н/м. Верхний брусок соединён невесомой нерастяжимой нитью, перекинутой через идеальный блок, с грузом массой $M = 2$~кг (см. рисунок). Чему равна длина пружины $l$ в нерастянутом состоянии, если при движении брусков её длина постоянна и равна $L = 15$~см? Сделайте рисунок с указанием сил, действующих на тела. Обоснуйте применимость используемых законов к решению задачи.

\nopagebreak\smallskip
\begin{center}
\begin{tikzpicture}[scale=1.2, every node/.style={font=\footnotesize}, decoration={coil, aspect=0.4, segment length=3pt, amplitude=2.5pt}]
    % Геометрия наклонной плоскости (угол примерно 26.5 градусов для удобства сетки: atan(1.5/3))
    % Координаты клина: (0,0) -> (3.5,0) -> (3.5,1.75) -> cycle
    \draw[thick, fill=gray!5] (0,0) -- (3.5,0) -- (3.5,1.75) -- cycle;
    \draw (0.5,0) arc (0:26.5:0.5);
    \node at (0.7,0.16) {$\alpha$};

    % Идеальный блок на вершине
    \coordinate (BlockCenter) at (3.5,1.75);
    \draw[thick, fill=white] (3.62,1.87) circle (0.15);
    \draw[thick, fill=white] (3.62,1.87) circle (0.04);
    \fill[black] (3.62,1.87) circle (1pt);
    \draw[thick] (BlockCenter) -- (3.62,1.87);

    % Наклонные линии для позиционирования брусков
    % Линия параллельна гипотенузе, смещена по нормали на высоту брусков
    % Нижний брусок m (первый по склону)
    \begin{scope}[shift={(1.0,0.5)}, rotate=26.5]
        \draw[thick, fill=brandAccent!20] (0,0) rectangle (0.5,0.3);
        \node at (0.25,0.45) {$m$};
    \end{scope}

    % Верхний брусок m (второй по склону)
    \begin{scope}[shift={(2.2,1.1)}, rotate=26.5]
        \draw[thick, fill=brandAccent!20] (0,0) rectangle (0.5,0.3);
        \node at (0.25,0.45) {$m$};
    \end{scope}

    % Пружина между брусками
    \begin{scope}[rotate=26.5]
        % В повернутых координатах торцы брусков находятся на фиксированных точках
        % Нижний брусок (от x=1.0 до x=1.5 в проекции) -> торец при x = 1.118*1.5 = 1.67
        \draw[thick, decorate] (1.56, 0.15) -- (2.32, 0.15);
    \end{scope}

    % Груз M (висит вертикально справа)
    \draw[thick, fill=brandPrimary!10] (3.62,0.4) rectangle (3.92,1.0);
    \node[right] at (3.95,0.7) {$M$};

    % Нити
    % Наклонная нить от верхнего бруска к блоку
    \begin{scope}[rotate=26.5]
        \draw[thick] (2.95,0.15) -- (3.9,0.15);
    \end{scope}
    % Вертикальная нить от блока к грузу M
    \draw[thick] (3.77,1.87) -- (3.77,1.0);

\end{tikzpicture}
\end{center}

\kim{26}
\end{taskbasic}
\answer{10~см}

% === ЗАДАЧА X ===
% Тег: #МЕХАНИКА #КИМ_26 #ДИНАМИКА_СВЯЗАННЫХ_ТЕЛ
\begin{taskbasic}[code={\label{task:elastic_friction_connected_n26}}]
\textbf{Задача X.} Груз массой $M = 800$~г соединён невесомой и нерастяжимой нитью, перекинутой через гладкий невесомый блок, с бруском массой $m = 400$~г. К этому бруску на лёгкой пружине жёсткостью $k = 80$~Н/м подвешен второй такой же брусок. Длина нерастянутой пружины $l = 10$~см, коэффициент трения груза о поверхность стола $\mu = 0{,}2$. Определите длину пружины при движении брусков, считая, что при этом движении она постоянна.

\nopagebreak\smallskip
\begin{center}
\begin{tikzpicture}[scale=1.2, every node/.style={font=\footnotesize}, decoration={coil, aspect=0.4, segment length=3pt, amplitude=2.5pt}]
    % Поверхность стола и обрыв стены (чистый fill=gray!5)
    \draw[thick, fill=gray!5] (-1.8,0) -- (0.2,0) -- (0.2,-0.1) -- (0.0,-0.1) -- (0.0,-2.5) -- (-0.1,-2.5) -- (-0.1,-0.1) -- cycle;
    % Штриховка стола и вертикальной стены
    \foreach \x in {-1.7,-1.5,-1.3,-1.1,-0.9,-0.7,-0.5,-0.3,-0.1}
        \draw[thin] (\x,0) -- (\x-0.1,-0.1);
    \foreach \y in {-0.2,-0.4,-0.6,-0.8,-1.0,-1.2,-1.4,-1.6,-1.8,-2.0,-2.2,-2.4}
        \draw[thin] (0.0,\y) -- (-0.1,\y-0.1);

    % Неподвижный блок на углу стола
    \draw[thick, fill=gray!30] (0.1,0) circle (0.14);
    \draw[thick, fill=gray!10] (0.1,0) circle (0.04);
    \fill[black] (0.1,0) circle (1pt);

    % Настольный груз M
    \draw[thick, fill=brandPrimary!10] (-1.2,0) rectangle (-0.5,0.6);
    \node at (-0.85,0.3) {$M$};

    % Верхний висящий брусок m
    \draw[thick, fill=brandAccent!20] (0.22,-0.9) rectangle (0.62,-0.5);
    \node at (0.82,-0.7) {$m$};

    % Пружина между брусками
    \draw[thick, decorate] (0.42,-0.9) -- (0.42,-1.6);

    % Нижний висящий брусок m
    \draw[thick, fill=brandAccent!20] (0.22,-2.0) rectangle (0.62,-1.6);
    \node at (0.82,-1.8) {$m$};

    % Горизонтальный участок нити (от M до блока)
    \draw[thick] (-0.5,0.14) -- (0.1,0.14);

    % Вертикальный участок нити (от блока до верхнего m)
    \draw[thick] (0.24,0) -- (0.24,-0.5);
\end{tikzpicture}
\end{center}

\kim{26}
\end{taskbasic}
\answer{13~см}

% === ЗАДАЧА X ===
% Тег: #МЕХАНИКА #КИМ_26 #МАШИНА_АТВУДА
\begin{taskbasic}[code={\label{task:atwood_overload_n26}}]
\textbf{Задача X.} Два одинаковых бруска массой $M = 500$~г связаны между собой невесомой нерастяжимой нитью, перекинутой через невесомый гладкий блок, неподвижно закреплённый на потолке (см. рисунок). На один из брусков кладут груз массой $m = 100$~г, и система приходит в движение. С какой силой $F$ груз будет давить на брусок? Сделайте схематический рисунок с указанием сил, действующих на бруски и груз. Обоснуйте применимость законов, используемых для решения задачи.

\nopagebreak\smallskip
\begin{center}
\begin{tikzpicture}[scale=1.2, every node/.style={font=\footnotesize}]
    % Потолок
    \draw[thick] (-0.8,2.5) -- (0.8,2.5);
    \foreach \x in {-0.7,-0.5,-0.3,-0.1,0.1,0.3,0.5,0.7}
        \draw[thin] (\x,2.5) -- (\x+0.1,2.6);

    % Крепление блока к потолку
    \draw[thick] (0,2.5) -- (0,1.8);

    % Неподвижный блок
    \draw[thick, fill=gray!10] (0,1.4) circle (0.4);
    \draw[thick, fill=white] (0,1.4) circle (0.08);
    \fill[black] (0,1.4) circle (1.2pt);

    % Левый брусок M
    \draw[thick, fill=brandPrimary!10] (-0.7,-0.4) rectangle (-0.1,0.1);
    \node[below] at (-0.4,-0.4) {$M$};

    % Правый брусок M
    \draw[thick, fill=brandPrimary!10] (0.1,-0.6) rectangle (0.7,-0.1);
    \node[below] at (0.4,-0.6) {$M$};

    % Маленький перегрузок m на правом бруске
    \draw[thick, fill=brandAccent!20] (0.25,-0.1) rectangle (0.55,0.2);
    \node[above] at (0.4,0.2) {$m$};

    % Нити от блока к брускам
    \draw[thick] (-0.4,1.4) -- (-0.4,0.1);
    \draw[thick] (0.4,1.4) -- (0.4,-0.1);

\end{tikzpicture}
\end{center}

\kim{26}
\end{taskbasic}
\answer{0{,}9~Н}

% === ЗАДАЧА X ===
% Тег: #МЕХАНИКА #КИМ_26 #ДВУМЕРНОЕ_ТРЕНИЕ
\begin{taskbasic}[code={\label{task:2d_friction_inclined_n26}}]
\textbf{Задача X.} На шероховатой наклонной плоскости, образующей с горизонтом угол $\alpha = 30^{\circ}$, лежит маленькая шайба массой $m = 500$~г. Минимальное значение модуля силы $\vec{F}$, приложенной в горизонтальном направлении вдоль плоскости (см. рисунок), при котором шайба начинает скользить по наклонной плоскости, равно $1{,}7$~Н. Чему равен коэффициент трения скольжения $\mu$ шайбы о плоскость?

\nopagebreak\smallskip
\begin{center}
\begin{tikzpicture}[scale=1.3, every node/.style={font=\footnotesize}]
    % Линии перспективы для наклонной плоскости
    % Основание (горизонтальный стол)
    \draw[thick] (0,0) -- (2.5,0);
    \draw[thick, dashed, gray] (2.5,0) -- (4.2,0.6);
    \draw[thick] (0,0) -- (1.7,0.6);
    
    % Задняя вертикальная грань клина
    \draw[thick] (4.2,0.6) -- (4.2,1.5);
    \draw[thick] (2.5,0) -- (4.2,0.6);
    
    % Наклонная поверхность (главный четырехугольник плоскости)
    \draw[thick, fill=gray!5] (0,0) -- (2.5,0) -- (4.2,1.5) -- (1.7,1.5) -- cycle;
    
    % Боковая правая грань клина для трехмерности
    \draw[thick, fill=gray!10] (2.5,0) -- (4.2,0.6) -- (4.2,1.5) -- cycle;

    % Обозначение угла альфа
    \draw (3,0.45) arc (45:10:0.4);
    \node at (3.4,0.5) {$\alpha$};

    % Шайба m (рисуем как маленький цилиндр в изометрии)
    \begin{scope}[shift={(1.8,0.75)}]
        % Нижнее основание цилиндра шайбы
        \draw[thick, fill=gray!40] (0,0) ellipse (0.22 and 0.09);
        \draw[thick, fill=gray!30](-0.22,0)--(-0.22,0.12);
        \draw[thick, fill=gray!30](0.22,0)--(0.22,0.12);
        % Стенки шайбы
       % \draw[thick, fill=gray!30] (-0.22,0) rectangle (0.22,0.12);
        % Верхнее основание цилиндра шайбы
        \draw[thick, fill=gray!50] (0,0.12) ellipse (0.22 and 0.09);
        
        % Точка на шайбе
        \fill[black] (0,0.06) circle (1pt);
        \node[left=4pt] at (-0.2,0.06) {$m$};
        
        % Вектор горизонтальной силы F (вдоль плоскости вбок)
        \draw[-{Stealth[scale=0.8]}, ultra thick, brandPrimary] (0,0.06) -- (0.8,0.06) node[above] {$\vec{F}$};
    \end{scope}

\end{tikzpicture}
\end{center}

\kim{26}
\end{taskbasic}
\answer{0{,}68}